\section{Einleitung}
\label{sec:einleitung}
In der heutigen digitalisierten Welt sind IT-Services ein integraler Bestandteil des Geschäftsbetriebs von nahezu allen Unternehmen, unabhängig von der Branche oder Größe.
Unternehmen sind auf IT-Services angewiesen, wie etwa grundlegende Netzwerkverwaltung, Cloud-Computing, fortschrittliche Cybersicherheit oder Entwicklung von Software.
Eine genaue Kalkulation und transparente Abrechnung dieser Services ist notwendig, um die erbrachten Leistungen kosteneffizient und fair verrechnen zu können.
\\
Die Berechnung und Abrechnung von IT-Dienstleistungen stellt jedoch eine komplexe Herausforderung dar.
Diese ergibt sich aus der Vielfalt der angebotenen Dienste, der variablen Nutzungsmuster sowie der unterschiedlichen Preismodelle, die von Dienstleistern verwendet werden.
\\
Diese Seminararbeit untersucht systematisch die verschiedenen Methoden und Ansätze zur Kalkulation und Abrechnung von IT-Services.
Ziel dieser Arbeit ist es, die wesentlichen Heausforderungen bei der Kalkulation und Abrechnung von IT-Services zu identifizieren und praxisnahe Lösungsansätze am Beispiel der K+s AG zu entwickeln.\\
Ein besonderes Augenmerk liegt auf der Definition und Strukturierung von IT-Services.
Dies ist notwendig, da diese die Basis für eine exakte Kalkulation und Abrechnung bilden.
Die Arbeit beleuchtet zudem eine notwendige kulturelle Transformation, um eine serviceorientierte Denkweise zu fördern, sowie die Umstellung der Kostenstrukturen von traditionellen Cost Centern zu Profit Centern.
\\
Durch die systematische Analyse und die Darstellung bewährter Praktiken sollen Unternehmen befähigt werden, effektive Strategien für die Kalkulation und Abrechnung ihrer IT-Services zu entwickeln und umzusetzen.
Dies soll dazu beitragen, die Transparenz und Effizienz von IT-Services zu steigern und den wirschaftlichen Erfolg des gesamten Unternehmens zu unterstützen.